%% ITEM 0 TYPE Section
\section{Introduction}%% ITEM 1 TYPE Frame
%% ITEM 1 TYPE Frame
\begin{frame}
\begin{center}
\Large \textbf{Introduction: LLM Post-Training Landscape}
\end{center}

\vspace{0.5cm}

\begin{itemize}
    \item \textbf{Two primary paradigms} dominate LLM post-training:
    \begin{itemize}
        \item \textbf{Supervised Fine-Tuning (SFT)}: Training on curated input-output pairs
        \item \textbf{Reinforcement Learning (RL)}: Optimizing against reward signals
    \end{itemize}
    
    \item \textbf{SFT} aligns models with human demonstrations through static datasets
    
    \item \textbf{RL} refines behavior in ways difficult to capture with SFT
    
    \item \textbf{RLHF} (Reinforcement Learning from Human Feedback) is a popular framework
    
    \item Algorithms like \textbf{PPO} and \textbf{GRPO} have become cornerstones
\end{itemize}
\end{frame}%% ITEM 2 TYPE Frame
%% ITEM 2 TYPE Frame
\begin{frame}
\begin{center}
\Large \textbf{The Need for a Unified View}
\end{center}

\vspace{0.5cm}

\begin{itemize}
    \item \textbf{Fragmented understanding}: SFT and RL are often treated as separate paradigms
    
    \item \textbf{Algorithm proliferation}: Many specialized algorithms (PPO, GRPO, CISPO, GSPO, etc.)
    
    \item \textbf{Limited theoretical connections}: Lack of common framework to understand relationships
    
    \item \textbf{Practical challenges}: Difficult to choose between approaches or combine them effectively
    
    \item \textbf{Research gap}: Need for systematic analysis of what makes each approach work
\end{itemize}
\end{frame}%% ITEM 3 TYPE Frame
%% ITEM 3 TYPE Frame
\begin{frame}
\begin{center}
\Large \textbf{Our Contribution: A Unified Framework}
\end{center}

\vspace{0.5cm}

\begin{itemize}
    \item \textbf{Unified perspective}: View both SFT and RL as post-training objectives
    
    \item \textbf{Common gradient form}: All algorithms can be expressed as:
    $$\text{grad}_{Uni} = \mathbb{1}_{stable} \frac{1}{\pi_{ref}} \hat{A} \nabla \pi_{\theta}$$
    
    \item \textbf{Four components}: Differences between algorithms break down into:
    \begin{itemize}
        \item Reference policy $\pi_{ref}$
        \item Advantage estimate $\hat{A}$
        \item Stability indicator $\mathbb{1}_{stable}$
        \item Gradient $\nabla \pi_{\theta}$
    \end{itemize}
    
    \item \textbf{Systematic analysis}: Enables comparison and understanding of algorithm relationships
\end{itemize}
\end{frame}%% ITEM 4 TYPE Frame
%% ITEM 4 TYPE Frame
\begin{frame}
\begin{center}
\Large \textbf{Summary and Roadmap}
\end{center}

\vspace{0.5cm}

\begin{itemize}
    \item \textbf{Current landscape}: SFT and RL as dominant post-training paradigms
    
    \item \textbf{Problem}: Fragmented understanding and algorithm proliferation
    
    \item \textbf{Solution}: Unified framework with common gradient representation
    
    \item \textbf{Key insight}: All algorithms share the same four-component structure
    
    \item \textbf{Next}: We'll explore the theoretical derivation and practical implications
\end{itemize}
\end{frame}%% ITEM 5 TYPE Section
%% ITEM 5 TYPE Section
\section{Unified Framework}%% ITEM 6 TYPE Frame
%% ITEM 6 TYPE Frame
\begin{frame}
\begin{center}
\Large \textbf{Unified Policy Gradient Estimator}
\end{center}

\vspace{0.3cm}

\textbf{The Common Objective:}
\begin{itemize}
    \item All post-training algorithms share a simple goal
    \item Maximize expected reward: $\max_{\theta} \mathcal{J}(\theta) := \mathbb{E}[r(\tau|q)]$
    \item Increase likelihood of positive trajectories
    \item Decrease likelihood of negative trajectories
    \item SFT and RL optimize the same fundamental objective
\end{itemize}
\end{frame}%% ITEM 7 TYPE Frame
%% ITEM 7 TYPE Frame
\begin{frame}
\textbf{The Unified Framework Components:}

\vspace{0.2cm}

Our gradient estimator decomposes into four key terms:

\begin{itemize}
    \item \textbf{Stabilization Mask}: Trust-region constraints (e.g., PPO clipping)
    \item \textbf{Reference Policy} $\pi_{ref}$: Sampling distribution for trajectories
    \item \textbf{Advantage Estimate} $\widehat{A}_{uni}$: Combines RL and SFT advantages
    \item \textbf{Likelihood Gradient} $\nabla_\theta \pi_\theta$: Policy gradient term
\end{itemize}

\vspace{0.2cm}

\textbf{Mathematical Form:}
\[
\mathrm{grad}_{uni} = \mathbb{1}_{stable} \cdot \frac{1}{\pi_{ref}} \cdot \hat{A} \cdot \nabla \pi_\theta
\]
\end{frame}%% ITEM 8 TYPE Frame
%% ITEM 8 TYPE Frame
\begin{frame}
\textbf{Unified Advantage Function:}

\vspace{0.2cm}

\[
\widehat{A}_{uni}(\tau,q) = \underbrace{r(\tau\mid q)}_{\widehat{A}_{\mathrm{RL}}(\tau,q)} \;+\; \underbrace{\mu\,\mathbb{1}\{\pi_{ref}=\pi_\beta\}\,\frac{\pi_\beta(\tau\mid q)}{\pi_\theta(\tau\mid q)}}_{\widehat{A}_{\mathrm{SFT}}(\tau,q)}
\]

\vspace{0.3cm}

\textbf{Key Insights:}
\begin{itemize}
    \item RL advantage: Raw reward $r(\tau|q)$ (often normalized)
    \item SFT advantage: Data adherence term with coefficient $\mu$
    \item Both terms can coexist in the same estimator
    \item Choice of $\pi_{ref}$ determines sampling distribution
    \item $\pi_{ref} = \pi_{\theta_{old}}$ for on-policy updates
    \item $\pi_{ref} = \pi_\beta$ for SFT/offline updates
\end{itemize}
\end{frame}%% ITEM 9 TYPE Frame
%% ITEM 9 TYPE Frame
\begin{frame}
\textbf{Full Gradient Expression:}

\vspace{0.2cm}

\[
\nabla_\theta \mathcal{J}_{\mu}(\theta) = \mathbb{E}_{\tau\sim \pi_{ref}(\cdot\mid q)}\!\left[
\frac{1}{\pi_{ref}(\tau\mid q)}\,
\widehat{A}_{uni}(\tau,q)\,
\nabla_\theta \pi_\theta(\tau\mid q)
\right]
\]

\vspace{0.3cm}

\textbf{Practical Implementation:}
\begin{itemize}
    \item With trust-region mask: $\mathbb{1}_{stable}(\tau,q)$ multiplies gradient
    \item RL works use structured advantages (e.g., GRPO normalization)
    \item Gradient is sum of two terms:
    \begin{itemize}
        \item Reward + trust-region term sampled from $\pi_\theta$
        \item Data-adherence (SFT) term sampled from $\pi_\beta$
    \end{itemize}
    \item Both map to same estimator via choice of $\pi_{ref}$
\end{itemize}
\end{frame}%% ITEM 10 TYPE Frame
%% ITEM 10 TYPE Frame
\begin{frame}
\textbf{Key Contributions \& Implications:}

\vspace{0.2cm}

\begin{itemize}
    \item \textbf{Unified View}: Shows SFT and RL optimize same objective
    \item \textbf{Modular Design}: Four components with clear roles
    \item \textbf{Joint Training}: SFT and RL can be trained together without conflict
    \item \textbf{Flexible Framework}: Accommodates various algorithms
    \begin{itemize}
        \item On-policy RL (PPO, GRPO)
        \item Offline RL
        \item Supervised fine-tuning
    \end{itemize}
    \item \textbf{Theoretical Foundation}: Derives from common reward maximization
    \item \textbf{Practical Impact}: Enables hybrid training approaches
\end{itemize}

\vspace{0.2cm}

\textbf{Conclusion}: A single framework that unifies seemingly different methods.
\end{frame}